\section{Implement Pow(x, n)}
\label{sec:rec-pow}

\problemstatement{Implement a function that computes $x^n$ for a real
number $x$ and an integer $n$, which may be negative, in better than
$O(|n|)$ time.}

\begin{patternbox}
\emph{``Compute $x^n$ efficiently''} is the canonical introduction to
\textbf{fast exponentiation} (also called binary exponentiation), already
used as a building block in Section~\ref{sec:rec-good-numbers}. The
keyword to internalize generally: whenever a recursive shrinking step can
\textbf{halve} the input rather than merely decrement it by one, the
recursion depth drops from $O(n)$ to $O(\log n)$ -- a pattern that
reappears constantly, not just for exponentiation.
\end{patternbox}

\subsection*{Understanding the problem carefully}

The naive recursive definition $x^n = x \times x^{n-1}$ (with base case
$x^0=1$) is correct but requires $n$ recursive calls. The key
algebraic identity that lets us do better is
\[
  x^n = \left(x^{n/2}\right)^2 \quad \text{when } n \text{ is even}, \qquad
  x^n = x \times \left(x^{(n-1)/2}\right)^2 \quad \text{when } n \text{ is
  odd}.
\]
Either way, computing $x^n$ only requires computing \emph{one} smaller
power, $x^{\lfloor n/2 \rfloor}$, and then a constant amount of extra
multiplication -- the exponent is halved at every recursive step instead
of merely decremented.

\begin{ideabox}[Halve the Exponent, Handle Negative Exponents Separately]
If $n=0$: return $1$ (base case). If $n<0$: compute $x^{-n}$ using the
same recursion, and return its reciprocal $1/x^{-n}$ at the very end (do
not try to handle negative exponents inside the recursive halving step
itself -- convert once, up front, and recurse only on non-negative
exponents). For $n>0$: recursively compute $\textit{half} =
x^{\lfloor n/2 \rfloor}$; if $n$ is even, return $\textit{half} \times
\textit{half}$; if $n$ is odd, return $\textit{half} \times \textit{half}
\times x$.
\end{ideabox}

\subsection*{Why halving gives logarithmic depth}

Each recursive call's exponent is at most half of its caller's exponent,
so after $k$ calls the exponent has shrunk by a factor of $2^k$ -- it
reaches the base case $0$ once $2^k \ge n$, i.e.\ after only
$O(\log n)$ calls, regardless of how large $n$ is. This is in sharp
contrast to the naive decrement-by-one recursion, whose depth is exactly
$n$. The trade-off is that each call here does a small constant amount
of extra multiplication work (squaring the smaller result, and possibly
one more multiplication by $x$ for odd exponents), but this constant
overhead is negligible compared to the exponential reduction in the
number of calls needed.

\subsection*{Algorithm}

\begin{enumerate}
  \item If $n<0$: return $1 / \textsc{Pow}(x, -n)$.
  \item $\textsc{Pow}(x, n)$:
  \begin{enumerate}
    \item If $n=0$: return $1$.
    \item $\textit{half} \gets \textsc{Pow}(x, n/2)$ (integer division).
    \item If $n$ is even: return $\textit{half} \times \textit{half}$.
    \item Else: return $\textit{half} \times \textit{half} \times x$.
  \end{enumerate}
\end{enumerate}

\subsection*{Dry run}

\dryrun{Compute $2^{10}$.}

\begin{itemize}
  \item $\textsc{Pow}(2,10)$: $10$ even. $\textit{half}=\textsc{Pow}(2,5)$.
  \item $\textsc{Pow}(2,5)$: $5$ odd. $\textit{half}=\textsc{Pow}(2,2)$.
  \item $\textsc{Pow}(2,2)$: $2$ even. $\textit{half}=\textsc{Pow}(2,1)$.
  \item $\textsc{Pow}(2,1)$: $1$ odd. $\textit{half}=\textsc{Pow}(2,0)=1$.
        Return $1\times1\times2=2$.
  \item Back at $\textsc{Pow}(2,2)$: $\textit{half}=2$. Return
        $2\times2=4$.
  \item Back at $\textsc{Pow}(2,5)$: $\textit{half}=4$. Return
        $4\times4\times2=32$.
  \item Back at $\textsc{Pow}(2,10)$: $\textit{half}=32$. Return
        $32\times32=1024$.
\end{itemize}

$2^{10}=1024$, confirmed. Only $4$ recursive calls were needed for an
exponent of $10$, instead of $10$ calls a naive decrement-by-one
recursion would have required.

\subsection*{C++ implementation}

\begin{lstlisting}
#include <bits/stdc++.h>
using namespace std;

double fastPow(double x, long long n) {
    if (n == 0) return 1.0;
    double half = fastPow(x, n / 2);
    double result = half * half;
    if (n % 2 == 1) result *= x;
    return result;
}

double myPow(double x, int n) {
    long long N = n;
    if (N < 0) return 1.0 / fastPow(x, -N);
    return fastPow(x, N);
}

int main() {
    cout << myPow(2.0, 10) << endl;   // 1024
    cout << myPow(2.0, -2) << endl;   // 0.25
    return 0;
}
\end{lstlisting}

\begin{complexitybox}
\textbf{Time Complexity.}
\[
  O(\log n).
\]
\textbf{Space Complexity.} $O(\log n)$ for the recursion stack (or
$O(1)$ if rewritten as an iterative loop that repeatedly squares the
base and halves the exponent, accumulating a result whenever the current
bit of the exponent is set).
\end{complexitybox}

\begin{takeawaybox}
Fast exponentiation is the prototype for an entire family of techniques:
whenever an operation is associative (like multiplication) and you need
to apply it $n$ times, check whether the operation can be ``doubled up''
(squared, in this case) to cut the required number of steps from $O(n)$
to $O(\log n)$. The same halving idea underlies fast matrix
exponentiation for linear recurrences and binary lifting for tree
ancestor queries, both common extensions of this exact pattern.
\end{takeawaybox}
